\section{Groups of Permutations}
\textbf{Conventions.} Permutations compose from right to left. The identity is $e$. A group generated by an element $a$ is denoted $\gen{a}$.

\exercise{8.1}{Product $\tau\sigma$}
\idea Apply the rightmost permutation first. The intermediate images are
\[\begin{array}{c|rrrrrr}x & 1&2&3&4&5&6\\\hline
\sigma(x) & 3 & 1 & 4 & 5 & 6 & 2\\
\tau(\sigma(x)) & 1 & 2 & 3 & 6 & 5 & 4\\\end{array}\]
The final row is the image of $1,\ldots,6$ under the required product. Thus\[\tau\sigma=\begin{pmatrix}1 & 2 & 3 & 4 & 5 & 6\\1 & 2 & 3 & 6 & 5 & 4\end{pmatrix}.\]

\exercise{8.2}{Product $\tau^2\sigma$}
\idea Apply the rightmost permutation first. The intermediate images are
\[\begin{array}{c|rrrrrr}x & 1&2&3&4&5&6\\\hline
\sigma(x) & 3 & 1 & 4 & 5 & 6 & 2\\
\tau(\sigma(x)) & 1 & 2 & 3 & 6 & 5 & 4\\
\tau(\tau(\sigma(x))) & 2 & 4 & 1 & 5 & 6 & 3\\\end{array}\]
The final row is the image of $1,\ldots,6$ under the required product. Thus\[\tau^2\sigma=\begin{pmatrix}1 & 2 & 3 & 4 & 5 & 6\\2 & 4 & 1 & 5 & 6 & 3\end{pmatrix}.\]

\exercise{8.3}{Product $\mu\sigma^2$}
\idea Apply the rightmost permutation first. The intermediate images are
\[\begin{array}{c|rrrrrr}x & 1&2&3&4&5&6\\\hline
\sigma(x) & 3 & 1 & 4 & 5 & 6 & 2\\
\sigma(\sigma(x)) & 4 & 3 & 5 & 6 & 2 & 1\\
\mu(\sigma(\sigma(x))) & 3 & 4 & 1 & 6 & 2 & 5\\\end{array}\]
The final row is the image of $1,\ldots,6$ under the required product. Thus\[\mu\sigma^2=\begin{pmatrix}1 & 2 & 3 & 4 & 5 & 6\\3 & 4 & 1 & 6 & 2 & 5\end{pmatrix}.\]

\exercise{8.4}{Product $\sigma^{-2}\tau$}
\idea Apply the rightmost permutation first. The intermediate images are
\[\begin{array}{c|rrrrrr}x & 1&2&3&4&5&6\\\hline
\tau(x) & 2 & 4 & 1 & 3 & 6 & 5\\
\sigma^{-1}(\tau(x)) & 6 & 3 & 2 & 1 & 5 & 4\\
\sigma^{-1}(\sigma^{-1}(\tau(x))) & 5 & 1 & 6 & 2 & 4 & 3\\\end{array}\]
The final row is the image of $1,\ldots,6$ under the required product. Thus\[\sigma^{-2}\tau=\begin{pmatrix}1 & 2 & 3 & 4 & 5 & 6\\5 & 1 & 6 & 2 & 4 & 3\end{pmatrix}.\]

\exercise{8.5}{Product $\sigma^{-1}\tau\sigma$}
\idea Apply the rightmost permutation first. The intermediate images are
\[\begin{array}{c|rrrrrr}x & 1&2&3&4&5&6\\\hline
\sigma(x) & 3 & 1 & 4 & 5 & 6 & 2\\
\tau(\sigma(x)) & 1 & 2 & 3 & 6 & 5 & 4\\
\sigma^{-1}(\tau(\sigma(x))) & 2 & 6 & 1 & 5 & 4 & 3\\\end{array}\]
The final row is the image of $1,\ldots,6$ under the required product. Thus\[\sigma^{-1}\tau\sigma=\begin{pmatrix}1 & 2 & 3 & 4 & 5 & 6\\2 & 6 & 1 & 5 & 4 & 3\end{pmatrix}.\]

\exercise{8.6}{Order of the subgroup generated by $\sigma$}
Following $1$ gives $1\mapsto3\mapsto4\mapsto5\mapsto6\mapsto2\mapsto1$. Thus $\sigma=(1,3,4,5,6,2)$. It takes exactly six applications to return $1$, and six applications fix every point. Therefore $\sigma$ has order $6$, so $|\gen{\sigma}|=6$.

\exercise{8.7}{Order of the subgroup generated by $	au^2$}
The cycles of $\tau$ are $(1,2,4,3)(5,6)$. Squaring advances two steps in each cycle, giving $\tau^2=(1,4)(2,3)$, with $5,6$ fixed. This permutation is not the identity, but its square is. Consequently $\gen{\tau^2}=\{e,\tau^2\}$ and the requested order is $2$.

\exercise{8.8}{The power $\sigma^{100}$}
Since $\sigma^6=e$ and $100=16\cdot6+4$, $\sigma^{100}=\sigma^4$. Advance four steps around $(1,3,4,5,6,2)$ to obtain
\[\sigma^{100}=\begin{pmatrix}1&2&3&4&5&6\\6&5&2&1&3&4\end{pmatrix}=(1,6,4)(2,5,3).\]
For example $1\mapsto3\mapsto4\mapsto5\mapsto6$, checking the first image.

\exercise{8.9}{The power $\mu^{100}$}
Here $\mu=(1,5)(3,4)$, so $\mu^2=e$. Since $100$ is even, $\mu^{100}=(\mu^2)^{50}=e$. The points $2,6$ are fixed throughout.

\exercise{8.10}{Partition into isomorphism classes}
\idea First separate finite groups by order, then use commutativity, torsion, and divisibility to distinguish infinite groups.
The complete partition is
\[\begin{gathered}
\{(\Z,+),(17\Z,+),(3\Z,+),\gen{\pi}\},\\
\{\Z_6,G\},\quad\{\Z_2,S_2\},\quad\{S_6\},\\
\{(\Q,+)\},\quad\{(\R,+),(\R^+,\cdot)\},\\
\{(\R^*,\cdot)\},\quad\{(\Q^*,\cdot)\},\quad\{(\C^*,\cdot)\}.
\end{gathered}\]
The maps $n\mapsto17n$, $n\mapsto3n$, and $n\mapsto\pi^n$ give the first class. The generator specified for $G$ is $(1,3,4)(2,5)$, of order $6$, so $G\cong\Z_6$. Both $S_2$ and $\Z_2$ are cyclic of order $2$. The group $S_6$ has $720$ elements and is nonabelian.
The exponential $x\mapsto e^x$ identifies $(\R,+)$ with $(\R^+,\cdot)$. The additive group $\Q$ is countable and divisible (every equation $nx=q$, $n>0$, has a solution), whereas $\Z$ is not divisible. The group $\Q^*$ is countable but has the element $-1$ of order $2$, distinguishing it from the countable torsion-free groups. The uncountable group $\R^*$ has exactly two torsion elements, $1,-1$, and is not divisible because $x^2=-1$ has no real solution. The group $\C^*$ has roots of unity of every positive order. These observations also exclude all remaining possible identifications.

\exercise{8.11}{Orbit under $\sigma$}
Repeatedly applying the permutation gives $1\mapsto3\mapsto4\mapsto5\mapsto6\mapsto2\mapsto1$. Negative powers traverse the same cycle backwards, so the orbit of $1$ is exactly $\{1,2,3,4,5,6\}$.

\exercise{8.12}{Orbit under $	au$}
Repeatedly applying the permutation gives $1\mapsto2\mapsto4\mapsto3\mapsto1$. Negative powers traverse the same cycle backwards, so the orbit of $1$ is exactly $\{1,2,3,4\}$.

\exercise{8.13}{Orbit under $\mu$}
Repeatedly applying the permutation gives $1\mapsto5\mapsto1$. Negative powers traverse the same cycle backwards, so the orbit of $1$ is exactly $\{1,5\}$.

\exercise{8.14}{Six geometric labels for $S_3$}
Let $r$ be a $120^\circ$ rotation of the equilateral triangle and $s$ its specified reflection. The rotations are $e,r,r^2$. Applying $s$ first and then each rotation gives the three reflections $s,rs,r^2s$, whose reflection axes pass through the three different vertices. They are distinct, and no reflection is a rotation because reflections reverse orientation. Hence $e,r,r^2,s,rs,r^2s$ are precisely the six elements of $S_3$.

\exercise{8.15}{Eight geometric labels for $D_4$}
Choose the quarter-turn $r=\rho_1$ and any one reflection $s$, for example $s=\mu_1$. The eight elements are
\[e,r,r^2,r^3,s,rs,r^2s,r^3s.\]
The first four preserve orientation and give the four rotational placements of a square. The last four reverse orientation and send a fixed vertex to four distinct vertices, so they are all different reflections. Thus they give every element of $D_4$.

\exercise{8.16}{Permutations fixing $3$}
Fixing $3$ in a permutation of $\{1,2,3,4\}$ leaves a free permutation of $\{1,2,4\}$. There are $3\cdot2\cdot1=6$ such choices, and each extends uniquely by $3\mapsto3$. The answer is $6$.

\exercise{8.17}{Permutations sending $2$ to $5$}
The image of $2$ is already fixed. The other four arguments must map bijectively to the four unused targets $1,2,3,4$. There are $4!=24$ ways to do this. Notice that fixing an image different from the argument does not change the count.

\exercise{8.18}{Subgroups of $S_3$}
\textbf{(a)} Both nonidentity rotations generate $\{e,\rho_1,\rho_2\}$, since $\rho_1^2=\rho_2$ and $\rho_1^3=e$. Since $\mu_1^2=e$, $\gen{\mu_1}=\{e,\mu_1\}$.
\textbf{(b)} The entire subgroup list is
\[\{e\},\quad\gen{\mu_1},\gen{\mu_2},\gen{\mu_3},\quad
\gen{\rho_1},\quad S_3.\]
An element is either a reflection or a rotation. Two distinct reflections have a nontrivial rotation as product, and a reflection together with a nontrivial rotation generates all six elements. Thus there are no additional subgroups. In the subgroup diagram $\{e\}$ lies below each of the four proper nontrivial subgroups, and each of those lies below $S_3$; there are no inclusions between those four.
\begin{center}\setlength{\unitlength}{1mm}\begin{picture}(115,36)
\put(52,31){$S_3$}\put(10,15){$\gen{\mu_1}$}\put(35,15){$\gen{\mu_2}$}\put(62,15){$\gen{\mu_3}$}\put(90,15){$\gen{\rho_1}$}\put(53,0){$\{e\}$}
\put(55,29){\line(-3,-1){39}}\put(55,29){\line(-1,-1){12}}\put(55,29){\line(1,-1){12}}\put(55,29){\line(3,-1){39}}
\put(55,5){\line(-3,1){39}}\put(55,5){\line(-1,1){12}}\put(55,5){\line(1,1){12}}\put(55,5){\line(3,1){39}}
\end{picture}\end{center}

\exercise{8.19}{All subgroups of $D_4$}
Write $r=\rho_1$ and $s=\mu_1$, so $r^4=s^2=e$ and $sr=r^{-1}s$. The cyclic subgroups generated by single elements are
\[\{e\},\quad\gen{r^2},\quad\gen{r},\quad\gen{s},\gen{rs},\gen{r^2s},\gen{r^3s}.\]
Here $\gen r$ has order $4$, while the five other nontrivial cyclic subgroups have order $2$. If a subgroup contains a reflection and $r$ or $r^3$, it is all of $D_4$. If it contains two reflections $r^is,r^js$, their product is $r^{i-j}$. For odd $i-j$ this again generates $D_4$; for even nonzero $i-j$, it produces one of the two Klein groups
\[K_0=\{e,r^2,s,r^2s\},\qquad K_1=\{e,r^2,rs,r^3s\}.\]
These are all possibilities for subgroups with two or more generators. Thus the diagram has top $D_4$; the next level $\gen r,K_0,K_1$; the next level the five order-$2$ groups; and bottom $\{e\}$. The group $\gen{r^2}$ connects to all three groups of order $4$, $\gen s,\gen{r^2s}$ connect to $K_0$, and $\gen{rs},\gen{r^3s}$ connect to $K_1$. These are exactly the edges in the supplied diagram, verifying it exhaustively.

\exercise{8.20}{The cyclic subgroup generated by $
ho$}
The permutation is $\rho=(1,2,4)(3,5)$, so its order is $\lcm(3,2)=6$. Indeed, $\rho^k=e$ precisely when both $3\mid k$ and $2\mid k$. In the following table $e=\rho^0=\rho^6$:
\[\begin{array}{c|cccccc}
 &e&\rho&\rho^2&\rho^3&\rho^4&\rho^5\\\hline
e & e & \rho^1 & \rho^2 & \rho^3 & \rho^4 & \rho^5\\
\rho^1 & \rho^1 & \rho^2 & \rho^3 & \rho^4 & \rho^5 & e\\
\rho^2 & \rho^2 & \rho^3 & \rho^4 & \rho^5 & e & \rho^1\\
\rho^3 & \rho^3 & \rho^4 & \rho^5 & e & \rho^1 & \rho^2\\
\rho^4 & \rho^4 & \rho^5 & e & \rho^1 & \rho^2 & \rho^3\\
\rho^5 & \rho^5 & e & \rho^1 & \rho^2 & \rho^3 & \rho^4\\\end{array}\]
The rule is $\rho^i\rho^j=\rho^{i+j\bmod6}$. This group is abelian and cyclic, so it is not isomorphic to the nonabelian group $S_3$, despite having the same order.

\exercise{8.21}{Six permutation matrices}
\textbf{(a)} Each displayed matrix has exactly one $1$ in each row and column, and all other entries zero. Such a matrix permutes the standard basis vectors. If $P_\sigma e_j=e_{\sigma(j)}$, then
\[P_\sigma P_\tau e_j=P_\sigma e_{\tau(j)}=e_{\sigma(\tau(j))},\]
so $P_\sigma P_\tau=P_{\sigma\tau}$. Thus the six matrices are closed under multiplication, contain $I$, and have inverses $P_{\sigma^{-1}}$; associativity comes from matrix multiplication.
\textbf{(b)} There are exactly $3!=6$ permutations of three basis vectors, and these are the six displayed matrices. The map $\sigma\mapsto P_\sigma$ is a bijective homomorphism, so the group is isomorphic to $S_3$.

\exercise{8.22}{Matrices representing $D_4$}
Use the rotation and reflection matrices
\[R=\begin{pmatrix}0&-1\\1&0\end{pmatrix},\qquad S=\begin{pmatrix}1&0\\0&-1\end{pmatrix}.\]
The requested eight matrices are $I,R,R^2,R^3,S,RS,R^2S,R^3S$, explicitly
\[\begin{gathered}
\begin{pmatrix}1&0\\0&1\end{pmatrix},\begin{pmatrix}0&-1\\1&0\end{pmatrix},
\begin{pmatrix}-1&0\\0&-1\end{pmatrix},\begin{pmatrix}0&1\\-1&0\end{pmatrix},\\
\begin{pmatrix}1&0\\0&-1\end{pmatrix},\begin{pmatrix}0&1\\1&0\end{pmatrix},
\begin{pmatrix}-1&0\\0&1\end{pmatrix},\begin{pmatrix}0&-1\\-1&0\end{pmatrix}.
\end{gathered}\]
They act on the square with vertices $(\pm1,\pm1)$ as exactly its eight symmetries. The equations $R^4=S^2=I$ and $SR=R^{-1}S$ prove closure and the same multiplication rule as $D_4$, giving the desired isomorphism.

\exercise{8.23}{Symmetries of the three-armed figure}
The three equal arms meet at $120^\circ$. Every symmetry fixes their common junction and permutes the three endpoints. Rotations through $0^\circ,120^\circ,240^\circ$ and the three reflections through an arm realize all six permutations of those endpoints. The action is faithful because an isometry fixing three noncollinear endpoints is the identity. Thus the group is $D_3\cong S_3$.

\exercise{8.24}{Symmetries of the divided circle}
The distinguished diameter must map to itself. The possibilities are the identity, the half-turn, reflection in that diameter, and reflection in its perpendicular diameter. Each nonidentity element squares to the identity, and the product of the two reflections is the half-turn. Hence the symmetry group is the Klein four-group $V\cong\Z_2\times\Z_2$.

\exercise{8.25}{Symmetries of the cross}
The four equal arms have endpoints forming a square. A symmetry of the cross gives a symmetry of those endpoints, and every symmetry of the square preserves the cross. The four rotations and four reflections therefore give exactly $D_4$.

\exercise{8.26}{Symmetries of the infinite decorated line}
The downward vertical teeth and upward slanting teeth distinguish the two sides of the horizontal line and also distinguish left from right. Thus neither a reflection nor a half-turn preserves the figure. The allowed translations are exactly integral multiples of the repeating horizontal spacing $a$. Their composition adds the integers: $T_{ma}T_{na}=T_{(m+n)a}$. Hence the symmetry group is the infinite cyclic group $\Z$.

\exercise{8.27}{Regular representations}
\textbf{Left regular representation of $\Z_4$.} For $k\in\Z_4$, put $L_k(x)=k+x\pmod4$. The four permutations are
\[L_0=e,\quad L_1=(0,1,2,3),\quad L_2=(0,2)(1,3),\quad L_3=(0,3,2,1),\]
with table
\[\begin{array}{c|cccc}&L_0&L_1&L_2&L_3\\\hline
L_0&L_0&L_1&L_2&L_3\\L_1&L_1&L_2&L_3&L_0\\
L_2&L_2&L_3&L_0&L_1\\L_3&L_3&L_0&L_1&L_2
\end{array}\]
\textbf{Right regular representation of $S_3$.} The book's homomorphic convention is $R_a(x)=xa^{-1}$, not $xa$. With columns in the order $\rho_0,\rho_1,\rho_2,\mu_1,\mu_2,\mu_3$, the images are
\[\begin{array}{c|cccccc}&\rho_0&\rho_1&\rho_2&\mu_1&\mu_2&\mu_3\\\hline
R_{\rho_0} & \rho_0 & \rho_1 & \rho_2 & \mu_1 & \mu_2 & \mu_3\\
R_{\rho_1} & \rho_2 & \rho_0 & \rho_1 & \mu_3 & \mu_1 & \mu_2\\
R_{\rho_2} & \rho_1 & \rho_2 & \rho_0 & \mu_2 & \mu_3 & \mu_1\\
R_{\mu_1} & \mu_1 & \mu_3 & \mu_2 & \rho_0 & \rho_2 & \rho_1\\
R_{\mu_2} & \mu_2 & \mu_1 & \mu_3 & \rho_1 & \rho_0 & \rho_2\\
R_{\mu_3} & \mu_3 & \mu_2 & \mu_1 & \rho_2 & \rho_1 & \rho_0\\\end{array}\]
For example $R_{\rho_1}(\rho_0)=\rho_2$. Since
$R_aR_b(x)=(xb^{-1})a^{-1}=x(ab)^{-1}$, the multiplication table of the representation is
\[\begin{array}{c|cccccc}&R_{\rho_0}&R_{\rho_1}&R_{\rho_2}&R_{\mu_1}&R_{\mu_2}&R_{\mu_3}\\\hline
R_{\rho_0} & R_{\rho_0} & R_{\rho_1} & R_{\rho_2} & R_{\mu_1} & R_{\mu_2} & R_{\mu_3}\\
R_{\rho_1} & R_{\rho_1} & R_{\rho_2} & R_{\rho_0} & R_{\mu_3} & R_{\mu_1} & R_{\mu_2}\\
R_{\rho_2} & R_{\rho_2} & R_{\rho_0} & R_{\rho_1} & R_{\mu_2} & R_{\mu_3} & R_{\mu_1}\\
R_{\mu_1} & R_{\mu_1} & R_{\mu_2} & R_{\mu_3} & R_{\rho_0} & R_{\rho_1} & R_{\rho_2}\\
R_{\mu_2} & R_{\mu_2} & R_{\mu_3} & R_{\mu_1} & R_{\rho_2} & R_{\rho_0} & R_{\rho_1}\\
R_{\mu_3} & R_{\mu_3} & R_{\mu_1} & R_{\mu_2} & R_{\rho_1} & R_{\rho_2} & R_{\rho_0}\\\end{array}\]
In each case evaluation at the identity distinguishes the permutations, so the maps from the original group are injective.

\exercise{8.28}{Definition of a permutation}
A permutation of $S$ is a \emph{bijection} $f:S\to S$: it must be both one-to-one and onto. The original wording omits surjectivity. For instance $n\mapsto n+1$ on $\{0,1,2,\ldots\}$ is one-to-one but misses $0$, and therefore is not a permutation.

\exercise{8.29}{Definition of the left regular representation}
The definition is correct: $g\mapsto L_g$, with $L_g(x)=gx$, maps $G$ into $S_G$. To check the terms are justified, $L_g^{-1}=L_{g^{-1}}$, so each $L_g$ is a permutation, and $L_gL_h=L_{gh}$. Finally $L_g(e)=g$, so this representation is one-to-one.

\exercise{8.30}{Translation as a permutation}
Yes. If $x+1=y+1$, subtracting $1$ gives $x=y$, so the map is one-to-one. Given a real target $u$, choosing $x=u-1$ gives $f(x)=u$, so it is onto. Its inverse is $u\mapsto u-1$.

\exercise{8.31}{Squaring as a permutation}
No. The two inputs $1$ and $-1$ have the same image $1$, so the map is not one-to-one. It also misses every negative real number because $x^2\ge0$. Either failure excludes it from $S_\R$.

\exercise{8.32}{The map $x\mapsto-x^3$}
Yes. If $-x^3=-y^3$, then $x^3=y^3$ and hence $x=y$, since the real cube function is strictly increasing. Given $u\in\R$, choose $x=\sqrt[3]{-u}$. Then $-x^3=u$, so the map is onto. Its inverse is $u\mapsto\sqrt[3]{-u}$.

\exercise{8.33}{The exponential map}
No. Although $e^x$ is strictly increasing and therefore one-to-one, its range is $(0,\infty)$, not all of $\R$. In particular there is no real input with image $0$. It would be a bijection from $\R$ to $\R^+$, but a permutation here must map $\R$ onto itself.

\exercise{8.34}{A cubic polynomial}
The polynomial is $f(x)=x^3-x^2-2x=x(x-2)(x+1)$. Thus $f(-1)=f(0)=f(2)=0$ at three distinct inputs, so it is not one-to-one and is not a permutation. (It is onto $\R$ by continuity and its opposite infinite end limits, but that does not repair injectivity.)

\exercise{8.35}{True or false}
\begin{enumerate}[label=(\alph*)]
\item True: a permutation is a bijection and hence injective.
\item False: injectivity does not ensure surjectivity or even equal domain and codomain; the shift on nonnegative integers is a counterexample.
\item True: if a finite set has $n$ elements, a surjection onto all $n$ targets cannot repeat an image, or at most $n-1$ targets would occur.
\item True: Cayley's theorem uses $g\mapsto L_g$ to embed $G$ into $S_G$.
\item True: if $ab=ba$ for all elements in a group, the same equality holds when both elements are restricted to any subgroup.
\item True: $\gen a=\{a^k:k\in\Z\}$ is closed under products and inverses and contains $e$.
\item False: $S_{10}$ has $10!$, not $10$, elements.
\item False: $S_3$ is nonabelian, whereas every cyclic group is abelian.
\item False: $S_1$ is trivial and cyclic, and $S_2$ is cyclic of order $2$.
\item True: this is the permutation-group form of Cayley's theorem.
\end{enumerate}

\exercise{8.36}{A nonabelian group with abelian proper subgroups}
Take $S_3$. Its proper subgroups are $\{e\}$, the three order-$2$ groups generated by transpositions, and the cyclic order-$3$ group $A_3$. All are cyclic and therefore abelian. Nevertheless $(1,2)(2,3)=(1,2,3)$ and $(2,3)(1,2)=(1,3,2)$ differ, so $S_3$ itself is nonabelian.

\exercise{8.37}{All self-maps under composition}
Function composition of maps $A\to A$ is an associative binary operation and has identity $\mathrm{id}_A$. Thus the set of all self-maps is a monoid (and in particular a semigroup). It is generally not a group: if $A$ has at least two elements, a constant map is not injective and cannot have an inverse. If $A$ is a singleton, there is only the identity map and it is the trivial group.

\exercise{8.38}{A Cayley digraph for $D_n$}
Let $r$ be a rotation of order $n$ and $s$ a reflection, so $sr=r^{-1}s$. Label the two rows (or two concentric $n$-gons) by $r^i$ and $r^is$, with $i$ read modulo $n$. For right-multiplication arcs, draw
\[r^i\xrightarrow{\ r\ }r^{i+1},\qquad
r^is\xrightarrow{\ r\ }r^{i-1}s,\qquad r^i\mathrel{\text{---}_{s}}r^is.\]
Thus the two $n$-cycles have opposite directions and a reflection edge joins corresponding vertices. This fully specifies the digraph, with $2n$ vertices and the correct outgoing arc of each type at every vertex. Schematically:
\[\begin{array}{cccccc}
e&\longrightarrow&r&\longrightarrow&r^2&\longrightarrow\cdots\longrightarrow r^{n-1}\longrightarrow e\\
\updownarrow s&&\updownarrow s&&\updownarrow s&\qquad\qquad\updownarrow s\\
s&\longleftarrow&rs&\longleftarrow&r^2s&\longleftarrow\cdots\longleftarrow r^{n-1}s\longleftarrow s
\end{array}\]
The reflection arcs are undirected because $s^2=e$.

\exercise{8.39}{Synopsis of Cayley's theorem}
Each $g\in G$ determines a permutation $L_g:x\mapsto gx$ of $G$, and $L_gL_h=L_{gh}$, so $g\mapsto L_g$ is a homomorphism into $S_G$. It is injective because $L_g(e)=g$, hence identifies $G$ with the subgroup $\{L_g:g\in G\}$.

\exercise{8.40}{The stabilizer of one point}
Yes. The identity fixes $b$. If $\sigma(b)=b$ and $\tau(b)=b$, then $(\sigma\tau)(b)=\sigma(b)=b$. Also $\sigma(b)=b$ implies $\sigma^{-1}(b)=b$. Thus the set contains the identity and is closed under composition and inverses, and is a subgroup of $S_A$.

\exercise{8.41}{Permutations sending a point into a subset}
Not necessarily. Let $A=\{1,2,3\}$, $B=\{1,2\}$, and $b=1$. The permutations $\sigma=(2,3)$ and $\tau=(1,2)$ both send $1$ into $B$, but $(\sigma\tau)(1)=\sigma(2)=3\notin B$. Thus the proposed set can fail closure under products.

\exercise{8.42}{Permutations mapping a subset into itself}
Not necessarily, because $B$ may be infinite. Let $A=\Z$ and $B=\{0,1,2,\ldots\}$. The shift $\sigma(n)=n+1$ satisfies $\sigma[B]\subseteq B$, but its inverse sends $0$ to $-1\notin B$. Hence the set can fail closure under inverses. If $B$ is finite, inclusion does force equality, but no finiteness is assumed here.

\exercise{8.43}{The setwise stabilizer of a subset}
Yes. The identity maps $B$ onto itself. If $\sigma[B]=B$ and $\tau[B]=B$, then $(\sigma\tau)[B]=\sigma[\tau[B]]=\sigma[B]=B$. Applying $\sigma^{-1}$ to $\sigma[B]=B$ gives $\sigma^{-1}[B]=B$. These three facts prove that the set is a subgroup.

\exercise{8.44}{Order and rotation subgroup of $D_n$}
Choose the image of one vertex: $n$ possibilities. Its adjacent vertex must go to one of the two adjacent vertices of the chosen image: $2$ possibilities. These choices determine the entire polygon placement, and all occur through rotations or reflections. Therefore $|D_n|=2n$. The rotations through multiples of $2\pi/n$ form the cyclic subgroup $\gen r=\{e,r,\ldots,r^{n-1}\}$ of order $n$, exactly half of $2n$.

\exercise{8.45}{Rotations of a cube}
Choose which of the six faces is on the bottom and then which of its four neighboring faces points forward. These choices determine a cube rotation, so the rotation group has $6\cdot4=24$ elements. For each of the three axes through centers of opposite faces, the quarter-turns form a distinct cyclic subgroup $\{e,r,r^2,r^3\}$ of order $4$. For each of the four long diagonals joining opposite vertices, the $120^\circ$ rotations form a distinct cyclic subgroup of order $3$. Different axes give different nonidentity rotations, establishing the requested lower bounds on numbers of subgroups.

\exercise{8.46}{Noncommutativity of $S_n$}
For $n\ge3$, let $a=(1,2)$ and $b=(2,3)$, fixing all other symbols. Then $(ab)(1)=a(1)=2$, whereas $(ba)(1)=b(2)=3$. Thus $ab\ne ba$ and $S_n$ is nonabelian. Comparing the images of one point is enough to distinguish two permutations.

\exercise{8.47}{Only the identity commutes with all permutations}
\idea Show a permutation moving a point fails to commute with a suitable transposition.
Suppose $\sigma(a)=b\ne a$. Choose $c$ different from $a,b$, possible because $n\ge3$, and let $\tau=(b,c)$. Then $\tau(a)=a$, so $(\sigma\tau)(a)=b$, while $(\tau\sigma)(a)=\tau(b)=c$. These images differ. Hence no nonidentity $\sigma$ commutes with every permutation. The identity does commute with every permutation, proving the claim.

\exercise{8.48}{Intersecting orbits coincide}
Suppose $c=\sigma^m(a)=\sigma^n(b)$ lies in both orbits. Applying $\sigma^{-n}$ gives $b=\sigma^{m-n}(a)$. Therefore any $\sigma^k(b)=\sigma^{k+m-n}(a)$ lies in the orbit of $a$, giving $O_{b,\sigma}\subseteq O_{a,\sigma}$. Interchanging $a,b$ gives the reverse inclusion, so the two orbits are equal.

\exercise{8.49}{A finite cyclic transitive permutation group}
Write the nonempty finite set as $A=\{a_1,\ldots,a_m\}$. Define $\sigma(a_i)=a_{i+1}$ for $i<m$ and $\sigma(a_m)=a_1$. The cycle $\sigma$ has order $m$, so $H=\gen\sigma$ has $|H|=m=|A|$. Given $a_i,a_j$, the power $\sigma^{j-i}$ carries $a_i$ to $a_j$ (with exponents modulo $m$). Thus $H$ is transitive. For $m=1$ the same construction gives the identity group.

\exercise{8.50}{Transitivity and a single orbit}
If $\gen\sigma$ is transitive and $a\in A$, then for every $b\in A$ some power $\sigma^k$ sends $a$ to $b$. Hence $O_{a,\sigma}=A$. Conversely, if this orbit is all of $A$, write any $b=\sigma^r(a)$ and $c=\sigma^s(a)$. Then $\sigma^{s-r}(b)=c$, so $\gen\sigma$ is transitive. The asserted existence of $a$ presupposes $A\ne\varnothing$; for the empty set transitivity is vacuous and this existential formulation needs that qualification.

\exercise{8.51}{The opposite group}
\textbf{(a)} Viewing the written equations through glass reverses left-to-right multiplication order: $a*b=c$ reads as $b*'a=c$. The associativity equation is reversed on both sides, so it is still associativity for the reversed operation. The identity and inverse equations are likewise preserved with the factors interchanged.
\textbf{(b)} Define $x*'y=y*x$. Closure follows from closure of $*$. Compute
\[(x*'y)*'z=z*(y*x)=(z*y)*x=x*'(y*'z),\]
using associativity of $*$. The old identity satisfies $e*'x=x*e=x$ and $x*'e=e*x=x$. The old inverse satisfies $x*'x^{-1}=x^{-1}*x=e$ and $x^{-1}*'x=x*x^{-1}=e$. Thus the reversed operation gives a group, called the opposite group.

\exercise{8.52}{The group of right translations}
For $a\in G$, define $\rho_a(x)=xa$. Each is a permutation with inverse $\rho_{a^{-1}}$. Composition reverses the subscripts:
\[\rho_a\rho_b(x)=(xb)a=x(ba),\qquad\rho_a\rho_b=\rho_{ba}.\]
Thus these permutations contain the identity, products, and inverses and form a group. The map $a\mapsto\rho_a$ need not be a homomorphism. Instead use $\Phi(a)=\rho_{a^{-1}}$; then
$\Phi(a)\Phi(b)=\rho_{b^{-1}a^{-1}}=\rho_{(ab)^{-1}}=\Phi(ab)$.
It is onto the set of right translations, and equality of two such translations, evaluated at $e$, gives $a^{-1}=b^{-1}$ and then $a=b$. Hence $\Phi$ is an isomorphism.

\exercise{8.53}{Permutation matrices and the Klein group}
\textbf{(a)} List the elements of a finite group $G$ in some order and let $L_g$ be left multiplication by $g$. Define $P_g e_x=e_{gx}$ on the standard basis indexed by $G$. Each $P_g$ is a permutation matrix. On every basis vector,
\[P_gP_h e_x=e_{ghx}=P_{gh}e_x,\]
so $g\mapsto P_g$ is a homomorphism. If $P_g=P_h$, applying both sides to $e_e$ gives $g=h$, so it is an isomorphism onto its image.
\textbf{(b)} Order $V$ as $e,a,b,c$, with $a^2=b^2=c^2=e$ and $ab=c$. Then
\[\begin{gathered}
P_e=I_4,\quad P_a=\begin{pmatrix}0&1&0&0\\1&0&0&0\\0&0&0&1\\0&0&1&0\end{pmatrix},\\
P_b=\begin{pmatrix}0&0&1&0\\0&0&0&1\\1&0&0&0\\0&1&0&0\end{pmatrix},\quad
P_c=\begin{pmatrix}0&0&0&1\\0&0&1&0\\0&1&0&0\\1&0&0&0\end{pmatrix}.
\end{gathered}\]
For example $P_a e_e=e_a$ and $P_a e_b=e_c$, matching multiplication by $a$. The construction in part (a) already verifies every product, including $P_aP_b=P_c$.
